\section{Conditional probability}
\label{conditionalProbabilitySection}

There can be rich relationships between two or more
variables that are useful to understand.
For example a car insurance company will consider
information about a person's driving history to assess
the risk that they will be responsible for an accident.
These types of relationships are the realm of conditional
probabilities.


\subsection{Exploring probabilities with a contingency table}

\index{data!photo\_classify|(}

\newcommand{\fashN}{1822}
% In order of ML, then Human
\newcommand{\fashYY}{197}
\newcommand{\fashYN}{22}
\newcommand{\fashYA}{219}
\newcommand{\fashNY}{112}
\newcommand{\fashNN}{1491}
\newcommand{\fashNA}{1603}
\newcommand{\fashAY}{309}
\newcommand{\fashAN}{1513}
\newcommand{\fashAA}{\fashN{}}
%\newcommand{\fashPYY}{}
%\newcommand{\fashPYN}{}
%\newcommand{\fashPNY}{}
%\newcommand{\fashPNN}{}
%\newcommand{\fashPYA}{0.12}
%\newcommand{\fashPNA}{0.88}
%\newcommand{\fashPAY}{}
%\newcommand{\fashPAN}{}
%\newcommand{\fashPYCY}{}
%\newcommand{\fashPYCN}{}
%\newcommand{\fashPNCY}{}
%\newcommand{\fashPNCN}{}
\newcommand{\fashCYPY}{0.96}
\newcommand{\fashCYPN}{0.04}
\newcommand{\fashCNPY}{0.07}
\newcommand{\fashCNPN}{0.93}

The \data{photo\us{}classify} data set represents
a classifier a sample of \fashN{} photos from a photo sharing website.
Data scientists have been working to improve a classifier for
whether the photo is about fashion or not, and these 1822 photos
represent a test for their classifier.
Each photo gets two classifications:
the first is called \var{mach\us{}learn} and gives
a classification from a machine
learning~(ML)\index{machine learning (ML)} system of
either \resp{pred\us{}fashion} or \resp{pred\us{}not}.
Each of these \fashN{} photos have also been classified carefully
by a team of people, which we take to be the source of truth;
this variable is called \var{truth} and takes values
\resp{fashion} and \resp{not}.
Figure~\ref{contTableOfFashionPhotos} summarizes the results.

\begin{figure}[ht]
\centering
\begin{tabular}{ll ccc rr}
&& \multicolumn{2}{c}{\var{truth}} & \hspace{1cm} &  \\
\cline{3-4}
&& \resp{fashion} & \resp{not} & Total  \\
\cline{2-5}
& \resp{pred\us{}fashion} &
    \fashYY{} & \fashYN{} & \fashYA{} \\
\raisebox{1.5ex}[0pt]{\var{mach\us{}learn}}
    & \resp{pred\us{}not} \hspace{0.5cm} &
    \fashNY{} & \fashNN{} & \fashNA{}   \\
\cline{2-5}
& Total & \fashAY{} & \fashAN{} & \fashN{} \\
\end{tabular}
\caption{Contingency table summarizing the
    \data{photo\us{}classify} data set.}
\label{contTableOfFashionPhotos}
\end{figure}
% library(openintro); table(photo_classify)

\begin{figure}[ht]
  \centering
  \Figure[A Venn diagram is shown, using boxes instead of circles, for the two categories of "ML Predicts Fashion" and "Fashion Photos" that partially overlap. The section of the rectangle for ML Predicts Fashion that is non-overlapping is labeled with 0.01. The section of the rectangle for Fashion Photos that is non-overlapping is labeled with 0.06. The overlapping section is labeled with 0.11. Outside of the rectangles is a label for "Neither" with a value 0.82.]{0.65}{photoClassifyVenn}
  \caption{A Venn diagram using boxes for the
      \data{photo\us{}classify} data set.}
  \label{photoClassifyVenn}
\end{figure}

\begin{examplewrap}
\begin{nexample}{If a photo is actually about fashion,
    what is the chance the ML classifier correctly identified
    the photo as being about fashion?}
  We can estimate this probability using the data.
  Of the \fashAY{} fashion photos,
  the ML algorithm correctly classified \fashYY{} of the photos:
\begin{align*}
P(\text{\var{mach\us{}learn} is \resp{pred\us{}fashion}
    given \var{truth} is \resp{fashion}})
  = \frac{\fashYY{}}{\fashAY{}}
  = 0.638
\end{align*}
\end{nexample}
\end{examplewrap}

\begin{examplewrap}
\begin{nexample}{We sample a photo from the data set
    and learn the ML algorithm predicted this photo
    was not about fashion.
    What is the probability that it was incorrect and
    the photo is about fashion?}
  If the ML classifier suggests a photo is not about fashion,
  then it comes from the second row in the data set.
  Of~these \fashNA{} photos, \fashNY{} were actually
  about fashion:
\begin{align*}
P(\text{\var{truth} is \resp{fashion}
    given \var{mach\us{}learn} is \resp{pred\us{}not}})
  = \frac{\fashNY{}}{\fashNA{}}
  = 0.070
\end{align*}
\end{nexample}
\end{examplewrap}

\subsection{Marginal and joint probabilities}
\label{marginalAndJointProbabilities}

\index{marginal probability|(}
\index{joint probability|(}

Figure~\ref{contTableOfFashionPhotos} includes row and
column totals for each variable separately in the
\data{photo\us{}classify} data set.
These totals represent
\termsub{marginal probabilities}{marginal probability}
for the sample, which are the probabilities based on a
single variable without regard to any other variables.
For instance, a probability based solely on the
\var{mach\us{}learn} variable is a marginal probability:
\begin{align*}
P(\text{\var{mach\us{}learn} is \resp{pred\us{}fashion}})
    = \frac{\fashYA{}}{\fashN{}}
    = 0.12
\end{align*}
A probability of outcomes for two or more variables
or processes is called a
\termsub{joint \mbox{probability}}{joint probability}:
\begin{align*}
P(\text{\var{mach\us{}learn} is \resp{pred\us{}fashion}
    and \var{truth} is \resp{fashion}})
  = \frac{\fashYY{}}{\fashN{}}
  = 0.11
\end{align*}
It is common to substitute a comma for ``and'' in a joint
probability, although using either the word ``and'' or a
comma is acceptable:
\begin{center}
$P(\text{\var{mach\us{}learn} is \resp{pred\us{}fashion},
    \var{truth} is \resp{fashion}})$ \\[2mm]
means the same thing as \\[2mm]
$P(\text{\var{mach\us{}learn} is \resp{pred\us{}fashion}
    and \var{truth} is \resp{fashion}})$
\end{center}

\begin{onebox}{Marginal and joint probabilities}
  If a probability is based on a single variable,
  it is a \emph{\hiddenterm{marginal probability}}.
  The probability of outcomes for two or more variables
  or processes is called a \emph{\hiddenterm{joint probability}}.
\end{onebox}

We use \term{table proportions} to summarize joint probabilities
for the \data{photo\us{}classify} sample.
These proportions are computed by dividing each count in
Figure~\ref{contTableOfFashionPhotos} by the table's total,
\fashN{}, to obtain the proportions in
Figure~\ref{photoClassifyProbTable}.
The joint probability distribution of the \var{mach\us{}learn}
and \var{truth} variables is shown in
Figure~\ref{photoClassifyDistribution}.

\begin{figure}[h]
\centering
\begin{tabular}{l rr r}
\hline
& \var{truth}: \resp{fashion} &
    \var{truth}: \resp{not} & Total  \\
\hline
\var{mach\us{}learn}: \resp{pred\us{}fashion} \hspace{0.5cm}
    & 0.1081 & 0.0121 & 0.1202 \\
\var{mach\us{}learn}: \resp{pred\us{}not}
    & 0.0615 & 0.8183 & 0.8798  \\
\hline
Total & 0.1696 & 0.8304 & 1.00 \\
\hline
\end{tabular}
\caption{Probability table summarizing the
    \var{photo\us{}classify} data set.}
\label{photoClassifyProbTable}
\end{figure}

\begin{figure}[h]
\centering
\begin{tabular}{l c}
  \hline
Joint outcome & Probability \\
  \hline
\var{mach\us{}learn} is \resp{pred\us{}fashion}
    and \var{truth} is \resp{fashion} & 0.1081 \\
\var{mach\us{}learn} is \resp{pred\us{}fashion}
    and \var{truth} is \resp{not} & 0.0121 \\
\var{mach\us{}learn} is \resp{pred\us{}not}
    and \var{truth} is \resp{fashion} & 0.0615 \\
\var{mach\us{}learn} is \resp{pred\us{}not}
    and \var{truth} is \resp{not} & 0.8183 \\
   \hline
Total & 1.0000 \\
\hline
\end{tabular}
\caption{Joint probability distribution for the \data{photo\us{}classify} data set.}
\label{photoClassifyDistribution}
\end{figure}

\D{\newpage}

\begin{exercisewrap}
\begin{nexercise}
Verify Figure~\ref{photoClassifyDistribution} represents
a probability distribution: events are disjoint,
all probabilities are non-negative, and the probabilities
sum to~1.\footnotemark
\end{nexercise}
\end{exercisewrap}
\footnotetext{Each of the four outcome combination are disjoint,
   all probabilities are indeed non-negative, and the sum of
   the probabilities is $0.1081 + 0.0121 + 0.0615 + 0.8183 = 1.00$.}

We can compute marginal probabilities using joint probabilities
in simple cases.
For example, the probability a randomly selected photo from the
data set is about fashion is found by summing the outcomes where
\var{truth} takes value \resp{fashion}:%
\index{marginal probability|)}\index{joint probability|)}
\newcommand{\ultruthfashion}[0]
    {\underline{\var{truth} is \resp{fashion}}}%
\begin{align*}
P(\text{\ultruthfashion{}})
  &= P(\text{\var{mach\us{}learn} is \resp{pred\us{}fashion}
        and \ultruthfashion{}}) \\
    & \qquad + P(\text{\var{mach\us{}learn} is \resp{pred\us{}not}
        and \ultruthfashion{}}) \\
  &= 0.1081 + 0.0615 \\
  &= 0.1696
\end{align*}


\subsection{Defining conditional probability}

\index{conditional probability|(}

The ML classifier predicts whether a photo is about fashion,
even if it is not perfect.
We would like to better understand how to use information
from a variable like \var{mach\us{}learn} to improve our
probability estimation of a second variable, which in this
example is \var{truth}.

The probability that a random photo from the data set is about
fashion is about 0.17.
If we knew the machine learning classifier predicted the
photo was about fashion, could we get a better estimate of the
probability the photo is actually about fashion?
Absolutely.
To do so, we limit our view to only those \fashYA{} cases
where the ML classifier predicted that the photo was about
fashion and look at the fraction where the photo was actually
about fashion:
\begin{align*}
P(\text{\var{truth} is \resp{fashion} given
    \var{mach\us{}learn} is \resp{pred\us{}fashion}})
  = \frac{\fashYY{}}{\fashYA{}}
  = 0.900
\end{align*}
We call this a \term{conditional probability} because
we computed the probability under a condition:
the ML classifier prediction said the photo was about fashion.

There are two parts to a conditional probability,
the \term{outcome of interest} and the \term{condition}.
It is useful to think of the condition as information we know
to be true, and this information usually can be described as
a known outcome or~event.
We generally separate the text inside our probability notation
into the outcome of interest and the condition with a
vertical bar:
\begin{align*}
&& P(\text{\var{truth} is \resp{fashion} given
    \var{mach\us{}learn} is \resp{pred\us{}fashion}}) \\
&& \quad = P(\text{\var{truth} is \resp{fashion}\ }|
    \text{\ \var{mach\us{}learn} is \resp{pred\us{}fashion}})
  = \frac{\fashYY{}}{\fashYA{}}
  = 0.900
\end{align*}
The vertical bar ``$|$'' is read as \emph{given}.

\D{\newpage}

In the last equation, we computed the probability a photo
was about fashion based on the condition that the ML algorithm
predicted it was about fashion as a fraction:
\begin{align*}
& P(\text{\var{truth} is \resp{fashion}\ }|
    \text{\ \var{mach\us{}learn} is \resp{pred\us{}fashion}}) \\
  &\quad = \frac{\text{\# cases where \var{truth} is \resp{fashion}
       and \var{mach\us{}learn} is \resp{pred\us{}fashion}}}
     {\text{\# cases where \var{mach\us{}learn} is \resp{pred\us{}fashion}}} \\
  &\quad = \frac{\fashYY{}}{\fashYA{}}
      = 0.900
\end{align*}
We considered only those cases that met the condition,
\var{mach\us{}learn} is \resp{pred\us{}fashion}, and then
we computed the ratio of those cases that satisfied our
outcome of interest, photo was actually about fashion.

Frequently, marginal and joint probabilities are provided
instead of count data.
For example, disease rates are commonly listed in percentages
rather than in a count format.
We would like to be able to compute conditional probabilities
even when no counts are available, and we use the last equation
as a template to understand this technique.

We considered only those cases that satisfied the condition,
where the ML algorithm predicted fashion.
Of these cases, the conditional probability was the
fraction representing the outcome of interest, that the
photo was about fashion.
Suppose we were provided only the information in
Figure~\ref{photoClassifyProbTable}, i.e. only probability data.
Then if we took a sample of 1000 photos, we would anticipate
about 12.0\% or $0.120\times 1000 = 120$ would be predicted to be
about fashion (\var{mach\us{}learn} is \resp{pred\us{}fashion}).
Similarly, we would expect about 10.8\% or
$0.108\times 1000 = 108$ to meet both the information criteria
and represent our outcome of interest.
Then the conditional probability can be computed as
\begin{align*}
&P(\text{\var{truth} is \resp{fashion}}\ |\ 
    \text{\var{mach\us{}learn} is \resp{pred\us{}fashion}}) \\
  &= \frac{\text{\# (\var{truth} is \resp{fashion}
      and \var{mach\us{}learn} is \resp{pred\us{}fashion})}}
    {\text{\# (\var{mach\us{}learn} is \resp{pred\us{}fashion})}} \\
  &= \frac{108}{120}
		= \frac{0.108}{0.120}
		= 0.90
\end{align*}
Here we are examining exactly the fraction of two probabilities,
0.108 and 0.120, which we can write as
\begin{align*}
P(\text{\var{truth} is \resp{fashion} and
    \var{mach\us{}learn} is \resp{pred\us{}fashion}})
\quad\text{and}\quad
P(\text{\var{mach\us{}learn} is \resp{pred\us{}fashion}}).
\end{align*}
The fraction of these probabilities is an example of the
general formula for conditional probability.

\begin{onebox}{Conditional probability}
  The conditional probability of outcome $A$
  given condition $B$ is computed as the following:
  \begin{align*}
  P(A | B) = \frac{P(A\text{ and }B)}{P(B)}
  \end{align*}
\end{onebox}

%\D{\newpage}

\begin{exercisewrap}
\begin{nexercise}
\label{fashionProbOfMLNotGivenTruthNot}%
(a) Write out the following statement in conditional
probability notation:
``\emph{The probability that the ML prediction was correct,
if the photo was about fashion}''.
Here the condition is now based on the photo's
\var{truth} status, not the ML algorithm. \\[1mm]
(b)~Determine the probability from part (a).
Table~\vref{photoClassifyProbTable} may be helpful.\footnotemark
\end{nexercise}
\end{exercisewrap}
\footnotetext{(a) If the photo is about fashion and the
  ML algorithm prediction was correct, then the ML algorithm
  my have a value of \resp{pred\us{}fashion}:
  \begin{align*}
  P(\text{\var{mach\us{}learn} is \resp{pred\us{}fashion}}\ |
      \ \text{\var{truth} is \resp{fashion}})
  \end{align*}
  (b)~The equation for conditional probability indicates we
  should first find \\
  $P(\text{\var{mach\us{}learn} is \resp{pred\us{}fashion}
    and \var{truth} is \resp{fashion}}) = 0.1081$
  and $P(\text{\var{truth} is \resp{fashion}}) = 0.1696$. \\
  Then the ratio represents the conditional probability:
  $0.1081 / 0.1696 = 0.6374$.}

\begin{exercisewrap}
\begin{nexercise}
\label{whyCondProbSumTo1}%
(a)~Determine the probability that the algorithm is incorrect
if it is known the photo is about fashion. \\[1mm]
(b)~Using the answers from part~(a) and
Guided Practice~\ref{fashionProbOfMLNotGivenTruthNot}(b),
compute
\begin{align*}
&P(\text{\var{mach\us{}learn} is \resp{pred\us{}fashion}}
    \ |\ \text{\var{truth} is \resp{fashion}}) \\
&\qquad  +\ P(\text{\var{mach\us{}learn} is \resp{pred\us{}not}}
    \ |\ \text{\var{truth} is \resp{fashion}})
\end{align*}
(c)~Provide an intuitive argument to explain why the sum
in~(b) is~1.\footnotemark
\end{nexercise}
\end{exercisewrap}
\footnotetext{(a)~This probability is
  $\frac{P(\text{\var{mach\us{}learn} is \resp{pred\us{}not},
      \var{truth} is \resp{fashion}})}
    {P(\text{\var{truth} is \resp{fashion}})}
  = \frac{0.0615}{0.1696} = 0.3626$.
  (b)~The total equals~1.
  (c)~Under the condition the photo is about fashion,
      the ML algorithm must have either predicted it was
      about fashion or predicted it was not about fashion.
      The complement still works for conditional probabilities,
      provided the probabilities are conditioned on the same
      information.}

\index{conditional probability|)}
\index{data!photo\_classify|)}


\subsection{Smallpox in Boston, 1721}

\index{data!smallpox|(}

The \data{smallpox} data set provides a sample of 6,224 individuals from the year 1721 who were exposed to smallpox in Boston.
Doctors at the time believed that inoculation, which involves exposing a person to the disease in a controlled form, could reduce the likelihood of death.

Each case represents one person with two variables: \var{inoculated} and \var{result}. The variable \var{inoculated} takes two levels: \resp{yes} or \resp{no}, indicating whether the person was inoculated or not. The variable \var{result} has outcomes \resp{lived} or \resp{died}. These data are summarized in Tables~\ref{smallpoxContingencyTable} and~\ref{smallpoxProbabilityTable}.

\begin{figure}[h]
\centering
\begin{tabular}{ll rr r}
& & \multicolumn{2}{c}{inoculated} & \\
\cline{3-4}
& & \resp{yes} & \resp{no} & Total  \\
\cline{2-5}
		& \resp{lived}     & 238 & 5136 & 5374 \\
\raisebox{1.5ex}[0pt]{\var{result}} &  \resp{died} \hspace{0.5cm} & 6 & 844 & 850  \\
\cline{2-5}
	& Total & 244 & 5980 & 6224 \\
\end{tabular}
\caption{Contingency table for the \data{smallpox} data set.}
\label{smallpoxContingencyTable}
\end{figure}

\begin{figure}[h]
\centering
\begin{tabular}{ll rr r}
& & \multicolumn{2}{c}{inoculated} & \\
\cline{3-4}
& & \resp{yes} & \resp{no} & Total  \\
   \cline{2-5}
 & \resp{lived}     & 0.0382 & 0.8252 & 0.8634 \\
\raisebox{1.5ex}[0pt]{\var{result}} & \resp{died} \hspace{0.5cm} & 0.0010 & 0.1356  & 0.1366  \\
   \cline{2-5}
& Total & 0.0392 & 0.9608 & 1.0000 \\
\end{tabular}
\caption{Table proportions for the \data{smallpox} data, computed by dividing each count by the table total, 6224.}
\label{smallpoxProbabilityTable}
\end{figure}

%\D{\newpage}

\begin{exercisewrap}
\begin{nexercise} \label{probDiedIfNotInoculated}
Write out, in formal notation, the probability a randomly selected person who was not inoculated died from smallpox, and find this \mbox{probability.}\footnotemark
\end{nexercise}
\end{exercisewrap}
\footnotetext{$P($\var{result} = \resp{died} $|$ \var{inoculated} = \resp{no}$) = \frac{P(\text{\var{result} = \resp{died} and \var{inoculated} = \resp{no}})}{P(\text{\var{inoculated} = \resp{no}})} = \frac{0.1356}{0.9608} = 0.1411$.}

\begin{exercisewrap}
\begin{nexercise}
Determine the probability that an inoculated person died from smallpox. How does this result compare with the result of Guided Practice~\ref{probDiedIfNotInoculated}?\footnotemark
\end{nexercise}
\end{exercisewrap}
\footnotetext{$P($\var{result} = \resp{died} $|$ \var{inoculated} = \resp{yes}$) = \frac{P(\text{\var{result} = \resp{died} and \var{inoculated} = \resp{yes}})}{P(\text{\var{inoculated} = \resp{yes}})} = \frac{0.0010}{0.0392} = 0.0255$ (if we avoided rounding errors, we'd get $6 / 244 = 0.0246$). The death rate for individuals who were inoculated is only about 1~in~40 while the death rate is about 1~in~7 for those who were not inoculated.}

\begin{exercisewrap}
\begin{nexercise}\label{SmallpoxInoculationObsExpExercise}
The people of Boston self-selected whether or not to be inoculated. (a) Is this study observational or was this an experiment? (b) Can we infer any causal connection using these data? (c) What are some potential confounding variables that might influence whether someone \resp{lived} or \resp{died} and also affect whether that person was inoculated?\footnotemark
\end{nexercise}
\end{exercisewrap}
\footnotetext{Brief answers: (a)~Observational. (b)~No, we cannot infer causation from this observational study. (c)~Accessibility to the latest and best medical care. There are other valid answers for part~(c).}


\subsection{General multiplication rule}

Section~\ref{probabilityIndependence} introduced the Multiplication Rule for independent processes. Here we provide the \term{General Multiplication Rule} for events that might not be independent.

\begin{onebox}{General Multiplication Rule}
If $A$ and $B$ represent two outcomes or events, then \vspace{-1.5mm}
\begin{align*}
P(A\text{ and }B) = P(A | B)\times P(B)
\end{align*} \vspace{-6.5mm} \par
It is useful to think of $A$ as the outcome of interest and $B$ as the condition.
\end{onebox}

\noindent%
This General Multiplication Rule is simply a rearrangement of the conditional probability equation.

%\D{\newpage}

\begin{examplewrap}
\begin{nexample}{Consider the \data{smallpox} data set. Suppose we are given only two pieces of information: 96.08\% of residents were not inoculated, and 85.88\% of the residents who were not inoculated ended up surviving. How could we compute the probability that a resident was not inoculated and lived?}
We will compute our answer using the General Multiplication Rule and then verify it using Figure~\ref{smallpoxProbabilityTable}. We want to determine
\begin{align*}
P(\text{\var{result}
    = \resp{lived} and \var{inoculated} = \resp{no}})
\end{align*}
and we are given that
\begin{align*}
P(\text{\var{result}
    = \resp{lived} }|\text{ \var{inoculated} = \resp{no}})
    &= 0.8588 %\\
&&P(\text{\var{inoculated} = \resp{no}})
    = 0.9608
\end{align*}
Among the 96.08\% of people who were not inoculated, 85.88\% survived:
\begin{align*}
P(\text{\var{result} = \resp{lived}
        and \var{inoculated} = \resp{no}})
    = 0.8588 \times 0.9608
    = 0.8251
\end{align*}
This is equivalent to the General Multiplication Rule. We can confirm this probability in Figure~\ref{smallpoxProbabilityTable} at the intersection of \resp{no} and \resp{lived} (with a small rounding error).
\end{nexample}
\end{examplewrap}

\begin{exercisewrap}
\begin{nexercise}
Use $P($\var{inoculated} = \resp{yes}$) = 0.0392$ and $P($\var{result} = \resp{lived} $|$ \var{inoculated} = \resp{yes}$) = 0.9754$ to determine the probability that a person was both inoculated and lived.\footnotemark
\end{nexercise}
\end{exercisewrap}
\footnotetext{The answer is 0.0382, which can be verified using Figure~\ref{smallpoxProbabilityTable}.}

%\D{\newpage}

\begin{exercisewrap}
\begin{nexercise}
If 97.54\% of the inoculated people lived,
what proportion of inoculated people must have died?\footnotemark{}
\end{nexercise}
\end{exercisewrap}
\footnotetext{There were only two possible outcomes: \resp{lived} or \resp{died}. This means that 100\% - 97.54\% = 2.46\% of the people who were inoculated died.}

\begin{onebox}{Sum of conditional probabilities}
Let $A_1$, ..., $A_k$ represent all the disjoint outcomes for a variable or process. Then if $B$ is an event, possibly for another variable or process, we have: \vspace{-1mm}
\begin{align*}
P(A_1|B) + \cdots + P(A_k|B) = 1
\end{align*}%
\vspace{-5.5mm} \par
The rule for complements also holds when an event and its complement are conditioned on the same information: \vspace{-1.5mm}
\begin{align*}
P(A | B) = 1 - P(A^c | B)
\end{align*}
\end{onebox}

\begin{exercisewrap}
\begin{nexercise}
Based on the probabilities computed above, does it appear that inoculation is effective at reducing the risk of death from smallpox?\footnotemark
\end{nexercise}
\end{exercisewrap}
\footnotetext{The samples are large relative to the difference in death rates for the ``inoculated'' and ``not inoculated'' groups, so it seems there is an association between \var{inoculated} and \var{outcome}. However, as noted in the solution to Guided Practice~\ref{SmallpoxInoculationObsExpExercise}, this is an observational study and we cannot be sure if there is a causal connection. (Further research has shown that inoculation is effective at reducing death rates.)}


%\D{\newpage}

\subsection{Independence considerations in conditional probability}

If two events are independent, then knowing the outcome of one should provide no information about the other. We can show this is mathematically true using conditional probabilities.

\begin{exercisewrap}
\begin{nexercise} \label{condProbOfRollingA1AfterOne1}
Let $X$ and $Y$ represent the outcomes of rolling two dice.\footnotemark
\begin{enumerate}[(a)]
\setlength{\itemsep}{0mm}
\item What is the probability that the first die, $X$, is \resp{1}?
\item What is the probability that both $X$ and $Y$ are \resp{1}?
\item Use the formula for conditional probability to compute $P(Y =$ \resp{1}$\ |\ X = $ \resp{1}$)$.
\item What is $P(Y=1)$? Is this different from the answer from part (c)? Explain.
\end{enumerate}
\end{nexercise}
\end{exercisewrap}
\footnotetext{Brief solutions: (a) $1/6$. (b) $1/36$. (c)~$\frac{P(Y = \text{ \resp{1} and }X=\text{ \resp{1}})}{P(X=\text{ \resp{1}})} = \frac{1/36}{1/6} = 1/6$. (d)~The probability is the same as in part~(c): $P(Y=1)=1/6$. The probability that $Y=1$ was unchanged by knowledge about $X$, which makes sense as $X$ and $Y$ are independent.}

We can show in Guided Practice~\ref{condProbOfRollingA1AfterOne1}(c) that the conditioning information has no influence by using the Multiplication Rule for independence processes:
\begin{align*}
P(Y=\text{\resp{1}}\ |\ X=\text{\resp{1}})
    &= \frac{P(Y=\text{\resp{1} and }X=\text{\resp{1}})}
      {P(X=\text{\resp{1}})} \\
    &= \frac{P(Y=\text{\resp{1}}) \times
        \color{oiGB}P(X=\text{\resp{1}})}
      {\color{oiGB}P(X=\text{\resp{1}})} \\
    &= P(Y=\text{\resp{1}}) \\
\end{align*}

\begin{exercisewrap}
\begin{nexercise}
Ron is watching a roulette table in a casino and notices that the last five outcomes were \resp{black}. He figures that the chances of getting \resp{black} six times in a row is very small (about $1/64$) and puts his paycheck on red. What is wrong with his reasoning?\footnotemark
\end{nexercise}
\end{exercisewrap}
\footnotetext{He has forgotten that the next roulette spin is independent of the previous spins. Casinos do employ this practice, posting the last several outcomes of many betting games to trick unsuspecting gamblers into believing the odds are in their favor. This is called the \term{gambler's fallacy}.}


\D{\newpage}

\subsection{Tree diagrams}

\index{data!smallpox|)}
\index{tree diagram|(}

\termsub{Tree diagrams}{tree diagram} are a tool to organize outcomes and probabilities around the structure of the data. They are most useful when two or more processes occur in a sequence and each process is conditioned on its predecessors.

The \data{smallpox} data fit this description. We see the population as split by \var{inoculation}: \resp{yes} and \resp{no}. Following this split, survival rates were observed for each group. This structure is reflected in the \term{tree diagram} shown in Figure~\ref{smallpoxTreeDiagram}. The first branch for \var{inoculation} is said to be the \term{primary} branch while the other branches are \termni{secondary}.

\begin{figure}[ht]
\centering
\Figure[A tree diagram with a primary branch "Inoculated" and a secondary branch "Result". The Inoculated primary branching leads to two options: "Yes" with a probability of 0.0392 and "No" with a probability of 0.9608. Each of these branches has secondary branches with conditional probabilities for the "Result" conditional on "Inoculated". The Inoculated Yes branch breaks into branches for "Lived" (0.9754) and "Died" (0.0246). These branches also provide the multiplied probabilities along the branches as well. For example, the Yes-and-Lived branching multiplies 0.0392 times 0.9754 to get 0.03824. The Yes-and-Died branching has a multiplied probability of 0.00096. Next, turning our attention to the "No" primary branch, it also has secondary branches of Lived and Died with conditional probabilities 0.8589 and 0.1411, respectively. It also shows the probabilities multiplied along each set of branches, with No-and-Lived as 0.82523 and No-and-Died as 0.13557.]{0.93}{smallpoxTreeDiagram}
\caption{A tree diagram of the \data{smallpox} data set.}
\label{smallpoxTreeDiagram}
\end{figure}

Tree diagrams are annotated with marginal and conditional probabilities, as shown in Figure~\ref{smallpoxTreeDiagram}. This tree diagram splits the smallpox data by \var{inoculation} into the \resp{yes} and \resp{no} groups with respective marginal probabilities 0.0392 and 0.9608. The secondary branches are conditioned on the first, so we assign conditional probabilities to these branches. For example, the top branch in Figure~\ref{smallpoxTreeDiagram} is the probability that \var{result} = \resp{lived} conditioned on the information that \var{inoculated} = \resp{yes}. We may (and usually do) construct joint probabilities at the end of each branch in our tree by multiplying the numbers we come across as we move from left to right. These joint probabilities are computed using the General Multiplication Rule:
\begin{align*}
& P(\text{\var{inoculated} = \resp{yes}
    and \var{result} = \resp{lived}}) \\
  &\quad = P(\text{\var{inoculated} = \resp{yes}})\times
      P(\text{\var{result} = \resp{lived}}|
          \text{\var{inoculated} = \resp{yes}}) \\
  &\quad = 0.0392\times 0.9754=0.0382
\end{align*}

\begin{examplewrap}
\begin{nexample}{Consider the midterm and final for a statistics class. Suppose 13\% of students earned an \resp{A} on the midterm. Of those students who earned an \resp{A} on the midterm, 47\% received an \resp{A} on the final, and 11\% of the students who earned lower than an \resp{A} on the midterm received an \resp{A} on the final. You randomly pick up a final exam and notice the student received an \resp{A}. What is the probability that this student earned an \resp{A} on the midterm?} \label{exerciseForTreeDiagramOfStudentGettingAOnMidtermGivenThatSheGotAOnFinal}
The end-goal is to find $P(\text{\var{midterm} = \resp{A}} | \text{\var{final} = \resp{A}})$. To calculate this conditional probability, we need the following probabilities:
\begin{align*}
P(\text{\var{midterm} = \resp{A} and \var{final} = \resp{A}})
  \qquad\text{and}\qquad
  P(\text{\var{final} = \resp{A}})
\end{align*}
However, this information is not provided, and it is not obvious how to calculate these probabilities. Since we aren't sure how to proceed, it is useful to organize the information into a tree diagram:
\begin{center}
\Figure[A tree diagram with a primary branch "Midterm" and a secondary branch "Final". The Midterm primary branching leads to two options: "A" with a probability of 0.13 and "Other" with a probability of 0.87. Each of these branches has secondary branches with conditional probabilities for the "Final" conditional on "Midterm". The Midterm-A branch breaks into branches for "A", with a conditional probability of 0.47 with an A-and-A final probability of 0.0611, and an "other" secondary branch, with a conditional probability of 0.53 with an other-and-other final probability of 0.0689. Next, turning our attention to the Midterm-Other primary branch, it also has secondary branches of Final-A with a conditional probability of 0.11 and final probability of 0.0957, and an "Final-other" branch with a conditional probability of 0.89 and final probability of 0.7743.]{0.85}{testTree}
\end{center}
When constructing a tree diagram, variables provided with marginal probabilities are often used to create the tree's primary branches; in this case, the marginal probabilities are provided for midterm grades. The final grades, which correspond to the conditional probabilities provided, will be shown on the secondary branches.

With the tree diagram constructed, we may compute the required probabilities:
\begin{align*}
&P(\text{\var{midterm} = \resp{A} and \var{final} = \resp{A}}) = 0.0611 \\
&P(\text{\underline{\color{black}\var{final} = \resp{A}}})  \\
& \quad= P(\text{\var{midterm} = \resp{other} and \underline{\color{black}\var{final} = \resp{A}}}) + P(\text{\var{midterm} = \resp{A} and \underline{\color{black}\var{final} = \resp{A}}}) \\
& \quad= 0.0957 + 0.0611  = 0.1568
\end{align*}
The marginal probability, $P($\var{final} = \resp{A}$)$, was calculated by adding up all the joint probabilities on the right side of the tree that correspond to \var{final} = \resp{A}. We may now finally take the ratio of the two probabilities:
\begin{align*}
P(\text{\var{midterm} = \resp{A}} | \text{\var{final} = \resp{A}})
  &= \frac{P(\text{\var{midterm} = \resp{A}
      and \var{final} = \resp{A}})}
    {P(\text{\var{final} = \resp{A}})} \\
  &= \frac{0.0611}{0.1568} = 0.3897
\end{align*}
The probability the student also earned an A on the midterm is about 0.39.
\end{nexample}
\end{examplewrap}

%\begin{figure}[ht]
%  \centering
%  \Figure{0.9}{testTree}
%  \caption{A tree diagram describing the \var{midterm}
%      and \var{final} variables.}
%  \label{testTree}
%\end{figure}

\begin{exercisewrap}
\begin{nexercise}
After an introductory statistics course, 78\% of students can successfully construct tree diagrams. Of those who can construct tree diagrams, 97\% passed, while only 57\% of those students who could not construct tree diagrams passed. (a)~Organize this information into a tree diagram. (b)~What is the probability that a randomly selected student passed? (c)~Compute the probability a student is able to construct a tree diagram if it is known that she passed.\footnotemark
\end{nexercise}
\end{exercisewrap}
\footnotetext{%\begin{minipage}[t]{0.27\linewidth}
(a) The tree diagram is shown to the right. \\
(b)~Identify which two joint probabilities represent
    students who passed, and add them:
    $P($passed$) = 0.7566+0.1254= 0.8820$. \\
(c)~$P($construct tree diagram $|$ passed$)
   = \frac{0.7566}{0.8820} = 0.8578$. \\ %\vspace{15mm} \\
%\end{minipage}
%\begin{minipage}[c]{0.7\linewidth}
\Figure[A tree diagram with a primary branch "Able to construct tree diagrams" and a secondary branch "Pass class". The Able-to-construct-tree-diagrams primary branching leads to two options: "Yes" with a probability of 0.78 and "No" with a probability of 0.22. Each of these branches has secondary branches with conditional probabilities for "Pass Class" conditional on "Able to construct tree diagrams". The Yes primary branch breaks into branches for "Pass", with a conditional probability of 0.97 with a Yes-and-Pass final probability of 0.7566, and a "Fail" secondary branch, with a conditional probability of 0.03 with a Yes-and-Fail final probability of 0.0234. Next, turning our attention to the No primary branch, it also has secondary branches of Pass with a conditional probability of 0.57 and final probability of 0.1254, and a Fail branch with a conditional probability of 0.43 and final probability of 0.0946.]{0.7}{treeDiagramAndPass}}% \vspace{-25mm}
%\end{minipage}}


\subsection{Bayes' Theorem}
\label{bayesTheoremSubsection}

\index{Bayes' Theorem|(}

In many instances, we are given a conditional probability of the form
\begin{align*}
P(\text{statement about variable 1 } | \text{ statement about variable 2})
\end{align*}
but we would really like to know the inverted conditional probability:
\begin{align*}
P(\text{statement about variable 2 } | \text{ statement about variable 1})
\end{align*}
Tree diagrams can be used to find the second conditional probability when given the first. However, sometimes it is not possible to draw the scenario in a tree diagram. In these cases, we can apply a very useful and general formula: Bayes' Theorem.

We first take a critical look at an example of inverting conditional probabilities where we still apply a tree diagram.

\D{\newpage}

\begin{examplewrap}
\begin{nexample}{In Canada, about 0.35\% of women over 40
    will develop breast cancer in any given year.
    A common screening test for cancer is the mammogram,
    but this test is not perfect.
    In about 11\% of patients with breast cancer, the test
    gives a \term{false negative}:
    it indicates a woman does not have breast cancer when
    she does have breast cancer.
    Similarly, the test gives a \term{false positive}
    in 7\% of patients who do not have breast cancer:
    it indicates these patients have breast cancer when
    they actually do not.
    If we tested a random woman over 40 for breast cancer
    using a mammogram and the test came back positive
    -- that is, the test suggested the patient has cancer --
    what is the probability that the patient actually has
    breast cancer?}

\label{probBreastCancerGivenPositiveTestExample}

Notice that we are given sufficient information to quickly compute the probability of testing positive if a woman has breast cancer ($1.00-0.11=0.89$). However, we seek the inverted probability of cancer given a positive test result. (Watch out for the non-intuitive medical language: a~\emph{positive} test result suggests the possible presence of cancer in a mammogram screening.) This inverted probability may be broken into two pieces:
\begin{align*}
P(\text{has BC } | \text{ mammogram$^+$}) = \frac{P(\text{has BC and mammogram$^+$})}{P(\text{mammogram$^+$})}
\end{align*}
where ``has BC'' is an abbreviation for the patient having
breast cancer and ``mammogram$^+$'' means the mammogram screening
was positive.
We can construct a tree diagram for these probabilities:
\begin{center}
\Figure[A tree diagram with a primary branch "Truth" and a secondary branch "Mammogram". The Truth primary branching leads to two options: "Cancer" with a probability of 0.0035 and "No Cancer" with a probability of 0.9965. Each of these branches has secondary branches with conditional probabilities for "Positive" and "Negative" mammogram outcomes conditional on whether the truth is having cancer or not. The Cancer primary branch breaks into branches for "Positive", with a conditional probability of 0.89 with a Cancer-and-Positive final probability of 0.00312, and a "Negative" secondary branch, with a conditional probability of 0.11 with a Cancer-and-Negative final probability of 0.00038. Next, turning our attention to the No-Cancer primary branch, it also has secondary branches of Positive with a conditional probability of 0.07 and final probability of 0.06976, and a Negative branch with a conditional probability of 0.93 and final probability of 0.92675.]{0.9}{BreastCancerTreeDiagram}
\end{center}
The probability the patient has breast cancer
and the mammogram is positive is
\begin{align*}
P(\text{has BC and mammogram$^+$}) &= P(\text{mammogram$^+$ } | \text{ has BC})P(\text{has BC}) \\
	&= 0.89\times 0.0035 = 0.00312
\end{align*}
The probability of a positive test result is the sum of the two corresponding scenarios:
\begin{align*}
P(\text{\underline{\color{black}mammogram$^+$}})
  &= P(\text{\underline{\color{black}mammogram$^+$} and has BC}) \\
  &\qquad\qquad
      + P(\text{\underline{\color{black}mammogram$^+$} and no BC})\\
  &= P(\text{has BC})P(\text{mammogram$^+$ } | \text{ has BC}) \\
  &\qquad\qquad
      + P(\text{no BC})P(\text{mammogram$^+$ } | \text{ no BC}) \\
  &= 0.0035\times 0.89 + 0.9965\times 0.07 = 0.07288
\end{align*}
Then if the mammogram screening is positive for a patient, the probability the patient has breast cancer is
\begin{align*}
P(\text{has BC } | \text{ mammogram$^+$})
	&= \frac{P(\text{has BC and mammogram$^+$})}{P(\text{mammogram$^+$})}\\
	&= \frac{0.00312}{0.07288} \approx 0.0428
\end{align*}
That is, even if a patient has a positive mammogram screening, there is still only a~4\%~chance that she has breast cancer.
\end{nexample}
\end{examplewrap}

%\begin{figure}[h]
%  \centering
%  \Figure{0.75}{BreastCancerTreeDiagram}
%  \caption{Tree diagram for
%      Example~\ref{probBreastCancerGivenPositiveTestExample}.}%,
%      %computing the probability a random patient who tests
%      %positive on a mammogram actually has breast cancer.}
%\label{BreastCancerTreeDiagram}
%\end{figure}

\D{\newpage}

Example~\ref{probBreastCancerGivenPositiveTestExample} highlights why doctors often run more tests regardless of a first positive test result. When a medical condition is rare, a single positive test isn't generally definitive.

Consider again the last equation of Example~\ref{probBreastCancerGivenPositiveTestExample}.
Using the tree diagram, we can see that the numerator (the top of the fraction) is equal to the following product:
\begin{align*}
P(\text{has BC and mammogram$^+$}) = P(\text{mammogram$^+$ } | \text{ has BC})P(\text{has BC})
\end{align*}
The denominator -- the probability the screening was positive -- is equal to the sum of probabilities for each positive screening scenario:
\begin{align*}
P(\text{\underline{\color{black}mammogram$^+$}})
	&= P(\text{\underline{\color{black}mammogram$^+$} and no BC})
		+ P(\text{\underline{\color{black}mammogram$^+$} and has BC})
\end{align*}
In the example, each of the probabilities on the right side was broken down into a product of a conditional probability and marginal probability using the tree diagram.
\begin{align*}
P(\text{mammogram$^+$})
	&= P(\text{mammogram$^+$ and no BC}) + P(\text{mammogram$^+$ and has BC}) \\
	&= P(\text{mammogram$^+$ } | \text{ no BC})P(\text{no BC}) \\
			   &\qquad\qquad + P(\text{mammogram$^+$ } | \text{ has BC})P(\text{has BC})
\end{align*}
We can see an application of Bayes' Theorem by substituting the resulting probability expressions into the numerator and denominator of the original conditional probability.
\begin{align*}
& P(\text{has BC } | \text{ mammogram$^+$})  \\
& \qquad= \frac{P(\text{mammogram$^+$ } | \text{ has BC})P(\text{has BC})}
	{P(\text{mammogram$^+$ } | \text{ no BC})P(\text{no BC}) + P(\text{mammogram$^+$ } | \text{ has BC})P(\text{has BC})}
\end{align*}

\begin{onebox}{Bayes' Theorem: inverting probabilities}
Consider the following conditional probability for variable 1 and variable 2:\vspace{-1.5mm}
\begin{align*}
P(\text{outcome $A_1$ of variable 1 } | \text{ outcome $B$ of variable 2})
\end{align*}
Bayes' Theorem states that this conditional probability can be identified as the following fraction:\vspace{-1.5mm}
\begin{align*}
\frac{P(B | A_1) P(A_1)}
	{P(B | A_1) P(A_1) + P(B | A_2) P(A_2) + \cdots + P(B | A_k) P(A_k)}
\end{align*}
where $A_2$, $A_3$, ..., and $A_k$ represent all other possible outcomes of the first variable.\index{Bayes' Theorem|textbf}
\end{onebox}

Bayes' Theorem is a generalization of what we have done
using tree diagrams.
The numerator identifies the probability of getting both
$A_1$ and~$B$.
The denominator is the marginal probability of getting~$B$.
This bottom component of the fraction appears long and
complicated since we have to add up probabilities from all of the different ways to get $B$. We always completed this step when using tree diagrams. However, we usually did it in a separate step so it didn't seem as complex.

\noindent%
To apply Bayes' Theorem correctly, there are two preparatory steps:
\begin{enumerate}
\setlength{\itemsep}{0mm}
\item[(1)] First identify the marginal probabilities of each possible outcome of the first variable: $P(A_1)$, $P(A_2)$, ..., $P(A_k)$.
\item[(2)] Then identify the probability of the outcome $B$, conditioned on each possible scenario for the first variable: $P(B | A_1)$, $P(B | A_2)$, ..., $P(B | A_k)$.
\end{enumerate}
Once each of these probabilities are identified, they can be applied directly within the formula.
Bayes' Theorem tends to be a good option when there are so many scenarios that drawing a tree diagram would be complex.

\begin{exercisewrap}
\begin{nexercise} \label{exerciseForParkingLotOnCampusBeingFullAndWhetherOrNotThereIsASportingEvent}
Jose visits campus every Thursday evening. However, some days the parking garage is full, often due to college events. There are academic events on 35\% of evenings, sporting events on 20\% of evenings, and no events on 45\% of evenings. When there is an academic event, the garage fills up about 25\% of the time, and it fills up 70\% of evenings with sporting events. On evenings when there are no events, it only fills up about 5\% of the time. If Jose comes to campus and finds the garage full, what is the probability that there is a sporting event? Use a tree diagram to solve this problem.\footnotemark
\end{nexercise}
\end{exercisewrap}
\footnotetext{\begin{minipage}[t]{0.27\linewidth}
The tree diagram, with three primary branches, is shown to the right. Next, we identify two probabilities from the tree diagram. (1) The probability that there is a sporting event and the garage is full: 0.14. (2) The probability the garage is full: $0.0875 + 0.14 + 0.0225 = 0.25$. Then the solution is the ratio of these probabilities: $\frac{0.14}{0.25} = 0.56$. If the garage is full, there is a 56\% probability that there is a sporting event. \vspace{0.1mm} \\\ 
\end{minipage}
\begin{minipage}[c]{0.65\linewidth}
\Figure[A tree diagram with a primary branch "Event" and a secondary branch "Garage full". The primary "Event" branching has three possibilities of "Academic" with probability 0.35, "Sporting" with probability 0.20, and "None" with probability 0.45. Each of these three branches has two secondary branches. The "Academic" primary branch breaks into branches for "Full" that has a conditional probability of 0.25 with an Academic-and-Full final probability of 0.0875, and a "Spaces Available" secondary branch with a conditional probability of 0.75 with an Academic-and-Spaces-Available final probability of 0.2625. The "Sporting" primary branch breaks into branches for "Full" that has a conditional probability of 0.7 with a Sporting-and-Full final probability of 0.14, and a "Spaces Available" secondary branch with a conditional probability of 0.3 with a Sporting-and-Spaces-Available final probability of 0.06. The "None" primary branch breaks into branches for "Full" that has a conditional probability of 0.05 with a None-and-Full final probability of 0.0225, and a "Spaces Available" secondary branch with a conditional probability of 0.95 with a None-and-Spaces-Available final probability of 0.4275.]{}{treeDiagramGarage}\vspace{-45mm}
\end{minipage}}

\begin{examplewrap}
\begin{nexample}{Here we solve the same problem presented in Guided Practice~\ref{exerciseForParkingLotOnCampusBeingFullAndWhetherOrNotThereIsASportingEvent}, except this time we use Bayes' Theorem.}
The outcome of interest is whether there is a sporting event (call this $A_1$), and the condition is that the lot is full ($B$). Let $A_2$ represent an academic event and $A_3$ represent there being no event on campus. Then the given probabilities can be written as
\begin{align*}
&P(A_1) = 0.2 &&P(A_2) = 0.35 &&P(A_3) = 0.45 \\
&P(B | A_1) = 0.7 &&P(B | A_2) = 0.25 &&P(B | A_3) = 0.05
\end{align*}
Bayes' Theorem can be used to compute the probability of a sporting event ($A_1$) under the condition that the parking lot is full ($B$):
\begin{align*}
P(A_1 | B) &= \frac{P(B | A_1) P(A_1)}{P(B | A_1) P(A_1) + P(B | A_2) P(A_2) + P(B | A_3) P(A_3)} \\
		&= \frac{(0.7)(0.2)}{(0.7)(0.2) + (0.25)(0.35) + (0.05)(0.45)} \\
		&= 0.56 
\end{align*}
Based on the information that the garage is full, there is a 56\% probability that a sporting event is being held on campus that evening.
\end{nexample}
\end{examplewrap}

\D{\newpage}

\begin{exercisewrap}
\begin{nexercise} \label{exerciseForParkingLotOnCampusBeingFullAndWhetherOrNotThereIsAnAcademicEvent}
Use the information in the previous exercise and example to verify the probability that there is an academic event conditioned on the parking lot being full is 0.35.\footnotemark
\end{nexercise}
\end{exercisewrap}
\footnotetext{Short answer:
\begin{align*}
P(A_2 | B) &= \frac{P(B | A_2) P(A_2)}{P(B | A_1) P(A_1) + P(B | A_2) P(A_2) + P(B | A_3) P(A_3)} \\
		&= \frac{(0.25)(0.35)}{(0.7)(0.2) + (0.25)(0.35) + (0.05)(0.45)} \\
		&= 0.35
\end{align*}}

\begin{exercisewrap}
\begin{nexercise} \label{exerciseForParkingLotOnCampusBeingFullAndWhetherOrNotThereIsNoEvent}
In Guided Practice~\ref{exerciseForParkingLotOnCampusBeingFullAndWhetherOrNotThereIsASportingEvent} and~\ref{exerciseForParkingLotOnCampusBeingFullAndWhetherOrNotThereIsAnAcademicEvent}, you found that if the parking lot is full, the probability there is a sporting event is 0.56 and the probability there is an academic event is 0.35. Using this information, compute $P($no event $|$ the lot is full$)$.\footnotemark
\end{nexercise}
\end{exercisewrap}
\footnotetext{Each probability is conditioned on the same information that the garage is full, so the complement may be used: $1.00 - 0.56 - 0.35 = 0.09$.}

The last several exercises offered a way to update our belief about whether there is a sporting event, academic event, or no event going on at the school based on the information that the parking lot was full. This strategy of \emph{updating beliefs} using Bayes' Theorem is actually the foundation of an entire section of statistics called \term{Bayesian statistics}. While Bayesian statistics is very important and useful, we will not have time to cover much more of it in this book.

\index{Bayes' Theorem|)}
\index{tree diagram|)}
\index{conditional probability|)}
\index{probability|)}


{\input{ch_probability/TeX/conditional_probability.tex}}
